\chapter{Conclusion and Future Work}

\section{Conclusion}
The main contribution of this thesis is not only the use of LLMs for task planning, but the
combination of four key elements in one multi-drone architecture: natural-language mission
interpretation, capability-aware allocation, explicit schedule generation, and local task admission
with replanning of unfinished subtasks. This distinguishes the proposed system from LLM-based
robotic planning approaches that generate sequential plans, dialogue-based coordination, or
reactive updates without explicit per-drone schedules and execution-time admission checking.

A cloud-based planning pipeline was developed in which a user mission is transformed into atomic subtasks,
assigned to suitable drones based on their capabilities, and scheduled while considering travel time and task duration.
The resulting output is an explicit per-drone schedule containing departure, arrival, and finish times.
In addition, a local task admission module was developed to evaluate whether individual drones are able to accept assigned
tasks based on their state, link quality, and available flight time.

The system was evaluated in three parts.
First, several cloud-based LLMs were compared on missions containing one to ten subtasks and evaluated against a
vehicle routing problem baseline. The results showed that GPT-5-mini provided the best trade-off between schedule
quality, validity, inference time, and cost, and was therefore selected for the cloud-based planner. Second, seven
quantized locally deployable models were evaluated for task admission. Among these, \texttt{qwen3:1.7b} achieved the
highest reliability and was selected for the admission module. Finally, the complete system was evaluated in a simulated
multi-drone environment, demonstrating that the generated plans could be translated into executable drone actions and
that unfinished or infeasible subtasks could be returned to the planner for rescheduling.

In summary, this thesis demonstrates that LLMs can serve as a high-level planning interface for
heterogeneous multi-drone systems when their outputs are constrained by structured task
representations, capability-aware allocation, explicit scheduling, and local feasibility checks.
The proposed architecture shows how natural-language mission planning can be connected to
schedule-based replanning, allowing unfinished subtasks to be reassigned when a drone rejects a
task or becomes unsuitable during execution.

\section{Limitations}
Although the system demonstrated promising results, several limitations should be considered when interpreting the findings.
The most important limitation is the inference time of the LLM-based planning pipeline. Long inference times limit the system's
suitability for time-critical missions and missions requiring frequent replanning. While smaller locally deployable models can
reduce latency, the results also showed that they are generally less reliable than larger cloud-based models. Similarly, the
inference time of the task admission module limits how frequently drone status checks can be performed during execution.

Another limitation is the possibility of invalid LLM outputs. Although structured prompts and validation mechanisms were used,
the planner may still occasionally produce malformed JSON outputs, missing fields, or schedules that violate system constraints.
Such invalid outputs can disrupt execution because they trigger replanning and thereby delay the mission.

In addition, the generated schedules do not guarantee optimality. As shown in the evaluation, LLM-generated schedules may be valid
and executable, but they are not always optimal with respect to the given optimization criterion. This is a limitation compared to
classical optimization methods, which can provide stronger guarantees under well-defined assumptions.

\section{Future Work}
Several directions can be explored to improve the proposed system. First, future work should investigate real-world
deployment on physical drones. While simulation provides a controlled environment for testing task planning and execution, 
real-world operation introduces additional challenges such as localization errors, communication delays, wind disturbances, 
obstacle avoidance, and sensor noise.

Second, the evaluation could be extended with a larger randomized benchmark. In this thesis, the cloud-based planner was
evaluated on missions containing one to ten subtasks, with three task descriptions per subtask count. This provides useful
initial coverage, but remains limited. A larger benchmark could include more mission descriptions, different object
distributions, different drone teams, and varied skill distributions. Such an evaluation would make it possible to assess 
more reliably whether the observed planning performance generalizes beyond the selected examples.

In addition, the task admission module could be extended with additional constraints, such as weather conditions, no-fly zones, 
payload limitations, obstacle proximity, and uncertainty in battery consumption. This would make the admission decision 
more realistic and better suited for deployment in dynamic environments.

Finally, future work could investigate more advanced local deployment of LLMs on embedded drone hardware. 
This includes optimizing models for hardware acceleration, reducing inference latency, and evaluating whether small onboard 
models can support not only task admission, but also limited replanning and failure recovery during flight.

In summary, this thesis demonstrates the potential of LLMs as a high-level interface between human operators and autonomous 
multi-drone systems. With further improvements in validation, real-world testing, and embedded deployment, LLM-based planning 
could become a useful component in future autonomous drone coordination systems.